%% ============================================================
%% 第六章 结论（中文版）
%% 结构：① 提出 QCCK + 作用 → ② QCCK 的内容 → ③ 进一步提出 QCCK-M
%%       → ④ 实验（最好的三个 + 可解释性） → ⑤ 意义与展望（意义写多一点）
%% 纪律：不写预测窗口、不写平均；实验只放 Beijing / ETTm2 / ECL 三个降幅。
%% 编译：xelatex（ctex）
%% ============================================================

\section{结论}

本文提出了量子跨变量耦合核 QCCK，用于在多变量时间序列预测中显式地构造变量之间的耦合关系。
与以扫描或注意力隐式传递变量信息的方式不同，QCCK 把变量间的依赖度量成一份可直接读取的矩阵，使跨变量建模的依据从黑盒权重中分离出来。

QCCK 由频域对齐、量子编码与量子核三个环节构成。
频域对齐在时间轴与频率轴上分别消除变量间的时移与周期尺度差异，把每个变量的原始序列变换为对齐频谱。
量子编码通过振幅通道与相位通道把变量自身特征与频域信息分别加载到量子态的振幅与相位上，得到纠缠量子态。
量子核以两个量子态的重叠程度度量变量间的耦合强度，输出取值在 0 到 1 之间的耦合矩阵 $K$。

在 QCCK 的基础上，本文进一步构建了多变量时间序列预测网络 QCCK-M。
QCCK-M 把 QCCK 以耦合层的形式嵌入主干编码器层之间，使耦合矩阵直接参与跨变量信息融合，在保留主干时序建模能力的同时引入显式的变量间依赖。

实验结果表明，QCCK-M 在多个公开数据集上优于现有方法，其中在 Beijing 数据集上 MSE 降低 11\%，在 ETTm2 上降低 6.2\%，在 ECL 上降低 6.1\%。
此外，耦合矩阵以显式数值呈现变量间的耦合强度，在 ETTh1 等系统上，模型读出的强耦合变量对与数据本身的物理分组一致，说明模型所依据的变量关系可以被直接读取与核对。

随着多变量时间序列应用场景的不断扩展，变量之间的关系往往是决定预测精度的关键，而这类关系通常难以从数据中直接读出。
QCCK 提供了一种把耦合显式化的建模途径，使变量间依赖既能被模型利用，也能被人核对；量子核在其中的作用表明，量子方法不仅可以用于提升表达能力，也可以为模型的可解释性提供支撑。
未来工作将进一步扩展耦合的刻画范围，使 QCCK 能够适应规模更大、耦合形态更复杂的多变量系统，并推动显式耦合建模在更多实际场景中的落地。
